% Mục 5: Phương pháp đo. Nguồn: scripts/run_bench.py (chú thích đầu file, khối "THE GRID"),
% include/bench_common.hpp, include/workload.hpp, scripts/plot.py, results/apple-m5/{calibration.json,
% campaign.log, *.csv}, tables.tex (campaign, lockunit). Số tự tính: scripts/report_numbers/corr_stats.py.
% Giới hạn của việc đo chỉ nằm ở mục Các yếu tố đe dọa độ tin cậy. Bản dài: agent-corr/long/method.tex.

Mỗi lần chạy in đúng một dòng CSV gồm tham số, thời gian, cột \texttt{ok} (phép kiểm của lần
chạy) và load average. Đồng hồ chạy từ lúc mọi thread
đã được tạo và cổng xuất phát chung mở ra tới khi lần \texttt{join} cuối trả về. Việc tạo thread nằm
ngoài vùng đo.

Trong khung của đề, compiler được phép xoá \texttt{is\_prime(N)} vì kết quả bị bỏ đi. Nhóm giữ
\texttt{is\_prime} không inline, cho \texttt{N} đi qua \texttt{pc::opaque()} ở mỗi lời gọi, và kiểm tổng các kết quả ở cuối lần chạy (cột \texttt{ok}).
Bảng~\ref{tab:apple-m5:lockunit} xác nhận phần việc không bị xoá. Ở $N=1009$, thời gian một lời gọi
$t$ chỉ thay đổi từ 504 tới 527~ns khi \texttt{OUTSIDE\_WORK} tăng từ 0 lên 15, và $t$ tỉ lệ với $N$.
Cả 12\,152 lần chạy đều có \texttt{ok = 1}.
% nguồn: tables.tex, bảng lockunit; campaign.log: bad_check=0

Phần lock quét 12 lock qua ba lượt. Một lượt dùng $N=1009$ với \texttt{OUTSIDE\_WORK} (OW) bằng 0, 1, 3, 7, 15. Hai lượt còn lại dùng $N=11$
và $N=10007$ với OW bằng 0 và 3. $N$ phải là số nguyên tố để \texttt{is\_prime} chạy đủ $N-2$ vòng. Ba
giá trị cho critical section khoảng 5~ns, 0,5~\textmu s và 5~\textmu s. OW quyết định mức tranh chấp.
Theo mô hình, lock bão hoà khi $p > 1+\text{OW}$. Barrier quét 7 cài đặt. Counter quét 19 cài đặt với OW bằng 0, 1, 8
và $N=101$. Số thread trên \PCname{apple-m5} là 1, 2, 3, 4, 6, 8, 10, 12, 16, 20. Trong đó 4 là số P-core, 10 là
mọi core, 12--20 là quá tải.

Với lock, số vòng mỗi thread là $I = W_N/\max(p,\,1+\text{OW})$, để cận dưới của mọi cấu hình khoảng
0,19~s. Nhóm hiệu chỉnh $W_N$ một lần cho mỗi $N$ bằng \PCname{ttas}. Với barrier và counter, chi phí
một thao tác chênh nhau nhiều bậc, nên mỗi cấu hình có $I$ riêng suy từ một lần chạy thử, nhắm 0,25~s
và 0,15~s. $I$ cố định cho cả đợt đo. Hình và bảng quy thời gian về cùng số vòng lặp.

Thời hạn mỗi lần chạy là 30~s cho lock và 10~s cho barrier và counter. Nhóm đo lock MCS quá tải với số
vòng nhỏ hơn 100 lần ngay từ đầu (27 cấu hình). Lần chạy vượt thời hạn được ghi thành một dòng
\texttt{ok=0}, không bỏ đi, rồi cấu hình đó chạy lại với số vòng nhỏ hơn 100 lần. Trên
\PCname{apple-m5} có 9 lần như vậy, đều là MCS quá tải (6 ở lock, 3 ở counter).
% nguồn: lock.csv, counter.csv (dòng extra=timeout); calibration.json (khoá reduced|...)

Mỗi cấu hình chạy 7 lần. Hình báo cáo trung vị, thanh sai số đi từ min tới max. Nhóm không dùng trung
bình vì phân bố lệch phải. Ở 33\,\% cấu hình lock, lần chậm nhất dài hơn 1,5 lần trung vị. Lần vượt
thời hạn được tính là chậm nhất khi lấy trung vị. Trong mỗi lượt lặp, mọi cấu hình chạy theo thứ tự xáo
trộn với seed cố định. Mỗi thời điểm chỉ
có một tiến trình đo, và đợt đo chỉ bắt đầu khi load average không quá 7,0.
% nguồn: agent-corr/stats.py trên lock.csv (ok=1): tỉ lệ cấu hình có max/median > 1,5 = 0,334

Phần lock được so với cận dưới $T_{\min} = \max\bigl(p\,I\,t,\ (1+\text{OW})\,I\,t\bigr)$. Vế đầu là
tổng các critical section chạy nối tiếp, vế sau là phần việc của riêng một thread. $t$ là trung vị lần
chạy 1 thread của lock nhanh nhất chia cho $I(1+\text{OW})$. Thời gian đo chia cho $T_{\min}$ là hệ số
chậm. Bảng~\ref{tab:apple-m5:campaign} tóm tắt đợt đo. Bảng~\ref{tab:apple-m5:lockunit} cho $t$.
